\documentclass{article}
\usepackage[T2A]{fontenc}
\usepackage[utf8]{inputenc}
\usepackage{amsthm}
\usepackage{amsmath}
\usepackage{amssymb}
\usepackage{amsfonts}
\usepackage{mathrsfs}
\usepackage[12pt]{extsizes}
\usepackage{fancyvrb}
\usepackage{indentfirst}
\usepackage[
  left=2cm, right=2cm, top=2cm, bottom=2cm, headsep=0.2cm, footskip=0.6cm, bindingoffset=0cm
]{geometry}
\usepackage[english,russian]{babel}


\begin{document}
\section*{Вариант 13}

Частные решения уравнений, соответствующие гармоническому закону пульсаций давления
на торце канала имеют вид:\

\begin{equation}
    R_k = \frac{2(-1)^{k+1}}{(2k-1) \pi} \frac{p^+_m}{w_mD}\left[A_{ck}\frac{df_p}{d\tau}+B_{ck}f_p\right],~~~
    Q_k = \frac{(-1)^{k+1}}{k \pi} \frac{p^+_m}{w_mD}\left[A_{sk}\frac{df_p}{d\tau}+B_{sk}f_p\right],
\end{equation}

где

\begin{equation*}
    A_{ck}=-\frac{a_{2ck}}{a^2_{1ck}+a^2_{2ck}},~~~
    B_{ck}=\frac{a_{1ck}}{a^2_{1ck}+a^2_{2ck}},~~~
    A_{sk}=-\frac{a_{2sk}}{a^2_{1sk}+a^2_{2sk}},~~~
    B_{sk}=\frac{a_{1sk}}{a^2_{1sk}+a^2_{2sk}}.

\end{equation*}

\end{document}
